% \documentclass[sigconf,edbt]{acmart-edbt2025}

% \def\BibTeX{{\rm B\kern-.05em{\sc i\kern-.025em b}\kern-.08em
%     T\kern-.1667em\lower.7ex\hbox{E}\kern-.125emX}}

% \usepackage{booktabs} % For formal tables


% % Copyright
% \setcopyright{rightsretained}

% % DOI
% \acmDOI{}

% % ISBN
% \acmISBN{978-3-89318-099-8}

% %Conference
% \acmConference[EDBT 2025]{28th International Conference on Extending Database Technology (EDBT)}{25th March-28th March, 2025}{Barcelona, Spain}
% \acmYear{2025}

% \settopmatter{printacmref=false, printccs=false, printfolios=false}

% \pagestyle{empty} % removes running headers

% \usepackage{graphicx}
% \usepackage[utf8]{inputenc}
% \usepackage{booktabs} % For formal tables
% \usepackage{fancyvrb}
% \usepackage{url}
% \usepackage{array}
% \usepackage{subcaption}
% \usepackage{amsfonts}
% \usepackage{amstext}
% \usepackage{amsmath}        % extensions for typesetting of math
  
% \usepackage[export]{adjustbox}
% \usepackage{multicol}
% \usepackage{multirow}
% \usepackage{algorithm2e}

% %%% This file contains definitions of various useful macros and environments %%%
%%% Please add more macros here instead of cluttering other files with them. %%%

\usepackage[HTML]{xcolor}   % typesetting in color
\usepackage{cleveref}       % better references

% The hyperref package for clickable links in PDF and also for storing
% metadata to PDF (including the table of contents).
% Most settings are pre-set by the pdfx package.
\hypersetup{
    unicode,
    breaklinks,
    colorlinks,
    linkcolor=blue,
    citecolor=blue!60!white,
    urlcolor=blue,
}

% Set up formatting of bibliography (references to literature)
% Details can be adjusted in macros.tex.
%
% BEWARE: Different fields of research and different university departments
% have their own customs regarding bibliography. Consult the bibliography
% format with your supervisor.
%
% The basic format according to the ISO 690 standard with numbered references
% \usepackage[natbib,style=iso-numeric,sorting=none]{biblatex}
% ISO 690 with alphanumeric references (abbreviations of authors' names)
%\usepackage[natbib,style=iso-alphabetic]{biblatex}
% ISO 690 with references Author (year)
%\usepackage[natbib,style=iso-authoryear]{biblatex}
%
% Some fields of research prefer a simple format with numbered references
% (sorting=none tells that bibliography should be listed in citation order)
\usepackage[natbib,style=numeric,sorting=none,urldate=long]{biblatex}
% Numbered references, but [1,2,3,4,5] is compressed to [1-5]
%\usepackage[natbib,style=numeric-comp,sorting=none]{biblatex}
% A simple format with alphanumeric references:
%\usepackage[natbib,style=alphabetic]{biblatex}

% Load the file with bibliography entries
\addbibresource{bibliography.bib}

%%% Switches based on thesis type

\def\TypeBc{bc}
\def\TypeMgr{mgr}
\def\TypePhD{phd}
\def\TypeRig{rig}

\ifx\ThesisType\TypeBc
\def\ThesisTypeName{bachelor}
\def\ThesisTypeTitle{BACHELOR THESIS}
\def\ThesisTypeCSGenitive{bakalářské}
\fi

\ifx\ThesisType\TypeMgr
\def\ThesisTypeName{master}
\def\ThesisTypeTitle{MASTER THESIS}
\def\ThesisTypeCSGenitive{diplomové}
\fi

\ifx\ThesisType\TypePhD
\def\ThesisTypeName{doctoral}
\def\ThesisTypeTitle{DOCTORAL THESIS}
\def\ThesisTypeCSGenitive{disertační}
\fi

\ifx\ThesisType\TypeRig
\def\ThesisTypeName{rigorosum}
\def\ThesisTypeTitle{RIGOROSUM THESIS}
\def\ThesisTypeCSGenitive{rigorozní}
\fi

\ifx\ThesisTypeName\undefined
\PackageError{thesis}{Unknown thesis type.}{Please check the definition of ThesisType in metadata.tex.}
\fi

%%% Switches based on study program language

\def\LangCS{cs}
\def\LangEN{en}

\ifx\StudyLanguage\LangCS
\else\ifx\StudyLanguage\LangEN
\else\PackageError{thesis}{Unknown study language.}{Please check the definition of StudyLanguage in metadata.tex.}
\fi\fi

%%% Minor tweaks of style

% These macros employ a little dirty trick to convince LaTeX to typeset
% chapter headings sanely, without lots of empty space above them.
% Feel free to ignore.
\makeatletter
\def\@makechapterhead#1{
  {\parindent \z@ \raggedright \normalfont
   \Huge\bfseries \thechapter\quad #1
   \par\nobreak
   \vskip 20\p@
}}
\def\@makeschapterhead#1{
  {\parindent \z@ \raggedright \normalfont
   \Huge\bfseries #1
   \par\nobreak
   \vskip 20\p@
}}
\makeatother

% This macro defines a chapter, which is not numbered, but is included
% in the table of contents.
\def\chapwithtoc#1{
\chapter*{#1}
\addcontentsline{toc}{chapter}{#1}
}

% Draw black "slugs" whenever a line overflows, so that we can spot it easily.
% \overfullrule=1mm

%%% Macros for definitions, theorems, claims, examples, ... (requires amsthm package)

% \theoremstyle{plain}
% \newtheorem{thm}{Theorem}
% \newtheorem{lemma}[thm]{Lemma}
% \newtheorem{claim}[thm]{Claim}
% \newtheorem{defn}{Definition}

% \theoremstyle{remark}
% \newtheorem*{cor}{Corollary}
% \newtheorem*{rem}{Remark}

\newtheoremstyle{examplestyle}% name
  {6pt}% Space above
  {6pt}% Space below
  {\small}% Body font
  {}% Indent amount
  {\bfseries}% Head font
  {.}% Punctuation after head
  {.5em}% Space after head
  {}% Head spec

\theoremstyle{examplestyle}
\newtheorem{exampleinner}{Example}[chapter]

\newenvironment{example}
  {\pushQED{\qed}\begin{exampleinner}}
  {\popQED\end{exampleinner}}

\crefname{exampleinner}{example}{examples}
\Crefname{exampleinner}{Example}{Examples}

%%% Style of captions of floating objects (figures etc.)

\ifcsname DeclareCaptionStyle\endcsname
\DeclareCaptionStyle{thesis}{style=base,font=small,labelfont=bf,labelsep=quad}
\captionsetup{style=thesis}
\captionsetup[algorithm]{style=thesis,singlelinecheck=off}
\captionsetup[listing]{style=thesis,singlelinecheck=off}
\fi

%%% An environment for typesetting of program code and input/output
%%% of programs.
% This would also define the `\code` command so fuck it.
% \DefineVerbatimEnvironment{code}{Verbatim}{fontsize=\small, frame=single}

% Settings for lstlisting -- program listing with syntax highlighting
\ifcsname lstset\endcsname
\lstset{
  % language=C++,
  tabsize=2,
  showstringspaces=false,
  basicstyle=\footnotesize\tt\color{black!75},
  identifierstyle=\bfseries\color{black},
  commentstyle=\color{green!50!black},
  stringstyle=\color{red!50!black},
  keywordstyle=\color{blue!75!black}}
\fi

% Floating listings, used in the same way as the figure environment
\ifcsname DeclareNewFloatType\endcsname
\DeclareNewFloatType{listing}{}
\floatsetup[listing]{style=ruled}
\floatname{listing}{Program}
\fi

%%% The field of all real and natural numbers
\newcommand{\R}{\mathbb{R}}
\newcommand{\N}{\mathbb{N}}

%%% Useful operators for statistics and probability
\DeclareMathOperator{\pr}{\textsf{P}}
\DeclareMathOperator{\E}{\textsf{E}}
\DeclareMathOperator{\var}{\textrm{var}}
\DeclareMathOperator{\sd}{\textrm{sd}}

%%% Transposition of a vector/matrix
\newcommand{\T}[1]{#1^\top}

%%% Asymptotic "O"
\def\O{\mathcal{O}}

%%% Various math goodies
\newcommand{\goto}{\rightarrow}
\newcommand{\gotop}{\stackrel{P}{\longrightarrow}}
\newcommand{\maon}[1]{o(n^{#1})}
\newcommand{\abs}[1]{\left|{#1}\right|}
\newcommand{\dint}{\int_0^\tau\!\!\int_0^\tau}
\newcommand{\isqr}[1]{\frac{1}{\sqrt{#1}}}

%%% TODO items: remove before submitting :)
\newcommand{\todo}[1]{\noindent{\color{red!}TODO: #1}}

%%% Detailed settings of bibliography

\ifx\citet\undefined\else

% Maximum number of authors of a single work. If exceeded, "et al." is used.
%\ExecuteBibliographyOptions{maxnames=2}
% The same setting specific to citations using \citet{...}
% \ExecuteBibliographyOptions{maxcitenames=2}
% The same settings specific to the list of literature
%\ExecuteBibliographyOptions{maxbibnames=2}

% Shortening first names of authors: "E. A. Poe" instead of "Edgar Allan Poe"
%\ExecuteBibliographyOptions{giveninits}
% The same without dots ("EA Poe")
%\ExecuteBibliographyOptions{terseinits}

% If your bibliography entries are hard to break into lines, try this mode:
%\ExecuteBibliographyOptions{block=ragged}

% Possibly reverse the names of the authors with the non-ISO styles:
%\DeclareNameAlias{default}{family-given}

% Use caps-and-small-caps for family names in ISO 690 style.
\let\familynameformat=\textsc

% We want to separate multiple authors in citations by commas
% (while we use semicolons in the bibliography as per the ISO standard)
\DeclareDelimFormat[textcite]{multinamedelim}{\addcomma\space}
\DeclareDelimFormat[textcite]{finalnamedelim}{\space and~}

\fi

\newcommand{\uv}[1]{``#1''}
\newcommand{\term}[1]{\textbf{#1}}
\newcommand{\code}[1]{{\normalfont{\texttt{#1}}}}
% Object from a schema category.
\newcommand{\object}[1]{\code{#1}}
% A signature of a morphism.
\newcommand{\signature}[1]{\code{#1}}

\makeatletter
\newcounter{Paper}

% Usage: \paperChapter[<TOC title>]{<label>}{<title>}{<authors>}{<publication>}
% Provide custom TOC title if the title is too long (it can overlap with the chapter number in the TOC).
\newcommand*{\paperChapter}[5][]{%
\refstepcounter{chapter}%
% This code is straight from latex2e source.
% Well, there is just `#3` instead of `\ifx...' in the source. However, this leads to overlapping of the title and the chapter number in the TOC, so we have to fix this with a custom TOC title.
\addcontentsline{toc}{chapter}{\protect\numberline{\thechapter}\ifx&#1&#3\else#1\fi}%
\refstepcounter{Paper}%
\edef\thisdisplay{\Roman{Paper}}%
\label{\detokenize{#2}}%
\protected@write\@auxout{}{%
  \string\newlabel{#2@display}{{\thisdisplay}{\thepage}}%
}%
\begin{center}
\hspace{0pt}
\vfill

{\LARGE\bfseries Paper~{\thisdisplay}\par}

\bigbreak

{\Large\bfseries #3\par}

\bigbreak\bigbreak

#4

\bigbreak\bigbreak

#5

\vfill
\hspace{0pt}
\end{center}
}

\newcount\paperRefCount

% Prints link to the given paper. Format: "I", "II", "III", ...
\newcommand{\paperRefNumber}[1]{\hyperref[\detokenize{#1}]{\ref{\detokenize{#1}@display}}}

% Prints reference to the given paper(s). Format: "Paper I", "Paper II", "Papers I, II, III". Only the number is hyperlinked.
\newcommand{\paperRef}[1]{%
  \edef\trimmedArguments{\zap@space#1 \@empty}%
  \paperRefCount=0\relax
  \@for\paper:=\trimmedArguments\do{\advance\paperRefCount by 1}%
  \ifnum\paperRefCount=1
    Paper~%
  \else
    Papers~%
  \fi
  \paperRefCount=0\relax
  \@for\paper:=\trimmedArguments\do{%
    \advance\paperRefCount by 1
    \ifnum\paperRefCount>1
      ,\penalty0\ %
    \fi
    \hyperref[\paper]{\ref{\paper @display}}%
  }%
}

\makeatother

\newcounter{autoFootnoteLabel}
\def\lastFootnoteLabel{}

\newcommand{\publishedAs}[1]{\noindent\textbf{Published as} \fullcite{#1}}

% Usage: \labeledFootnote[<label>]{<text>}
% Creates a footnote with a specific label. If the label is not provided, it will be generated automatically (a number).
\NewDocumentCommand{\labeledFootnote}{O{} m}{%
  % Determine label
  \def\currentLabel{#1}%
  \ifx\currentLabel\empty
    \stepcounter{autoFootnoteLabel}%
    \edef\currentLabel{autofn:\arabic{autoFootnoteLabel}}%
  \fi
  % Remember for later \footref
  \global\let\lastFootnoteLabel\currentLabel
  % Create footnote
  \footnote{\label{\currentLabel}#2}%
}

% Usage: \footnoteRef[<label>]
% Creates a reference to a footnote with the given label. If the label is not provided, it will refer to the last created `\labeledFootnote`.
\NewDocumentCommand{\footnoteRef}{O{}}{%
  \def\currentLabel{#1}%
  \ifx\currentLabel\empty
    \let\currentLabel\lastFootnoteLabel
  \fi
  \hyperref[\currentLabel]{\footnotemark[\getrefnumber{\currentLabel}]}%
}

% Common stuff for all papers
\newcommand{\mffAuthor}{Department of Software Engineering, Faculty of Mathematics and Physics, Charles University, Prague, Czech Republic}

\newcommand{\footnotei}[2]{%
\mbox{%
\setbox0\hbox{#2}%
\copy0%
\hspace{-\wd0}}%
\footnote{#1}%
}

% Width to use for a figure in a fbox to make it fit the page without overflowing. Measured experimentally to a sufficient precision.
\def\framedWidth{0.999275\textwidth}

% Decides whether to include the papers. See metadata.tex for the actual switch.
% Overrides labels and refs for papers.
\def\true{true}

% Usage: \paperInput{<label>}{<path>}
% Includes a paper from the given path. All \pLabel and \pRef commands in the paper will be prefixed with the given label.
\newcommand{\paperInput}[2]{%
\ifx\includePapers\true
\newpage
\begingroup
  \def\inputPrefix{#1}%

  % This approach doesn't work for \label, because it's being redefined by someone (probably at sections or sth). Yup, macros at its finest.
  % So, we have to use the \pLabel macro below.
  % \let\oldlabel\label
  % \renewcommand{\label}[1]{\oldlabel{\inputPrefix##1}}

  % So, this doesn't work either ... gues what, if \ref is used in a caption, it's written expdanded to \oldRef in the .aux file, which is, ofc, undefined when used in LOF.
  % \let\oldRef\ref
  % \renewcommand{\ref}[1]{\oldRef{\inputPrefix##1}}

  \input{#2}
\endgroup
\fi
}

% Use this instead of \label and \ref in the included papers.
\newcommand{\pLabel}[1]{\label{\inputPrefix#1}}
\newcommand{\pRef}[1]{\ref{\inputPrefix#1}}

% Usage: \defineQuestionPrefix{<prefix>}{<marker>}
% Defines a new marker for the given question prefix. E.g., if the marker is "Q", the questions will be numbered as "Q1", "Q2", etc.
\newcommand{\defineQuestionPrefix}[2]{%
  \newcounter{question@#1}%
  \expandafter\def\csname questionprefix@#1\endcsname{#2}%
}

% Usage: \question{<prefix>}{<label>}{<human-readable label>}{<question text>}
% Creates a question
\newcommand{\question}[4]{%
  \refstepcounter{question@#1}%
  \par\medskip
  \noindent\textbf{\csname questionprefix@#1\endcsname\arabic{question@#1}}\label{#1:#2} \ \emph{#3} \ #4\par\medskip
}

% Usage: \questionRef{<prefix>}{<label>}
% References a question
\newcommand{\questionRef}[2]{%
  \hyperref[#1:#2]{%
    \csname questionprefix@#1\endcsname\ref{#1:#2}%
  }%
}

\newcounter{publication}
\newcommand{\publication}[1]{%
\refstepcounter{publication}%
\noindent\makebox[2em][l]{\textbf{\arabic{publication}.}}%
\hangindent=2em\hangafter=1\relax%
\fullcite{#1}\par
\vspace{1em}
}

% Settings for the alltt environment (like, they are global, but we use them only there).

\newcommand{\keyword}[1]{\textbf{#1}}

\definecolor{sql}{HTML}{9673A6}
\newcommand{\sql}[1]{\textcolor{sql}{\bfseries #1}}

\definecolor{cypher}{HTML}{6C8EBF}
\newcommand{\cypher}[1]{\textcolor{cypher}{\bfseries #1}}

\definecolor{xquery}{HTML}{82B366}
\newcommand{\xquery}[1]{\textcolor{xquery}{\bfseries #1}}

\newcounter{contribution}[section]
\newcommand{\contribution}[1]{%
  \refstepcounter{contribution}%
  \paragraph{\arabic{contribution}}\emph{#1} \ %
}
  % \par\medskip
  % \noindent\textbf{\arabic{contribution}} \ \emph{#1} \ #2\par\medskip


% \begin{document}
% \title{\textit{TransforMMer}: A Universal Multi-Model Data Generator}

% \author{Jáchym Bártík}
% \orcid{0000-0002-5664-5890}
% \affiliation{%
%   \institution{Charles University}
%   \city{Prague} 
%   \country{Czech Republic}
% }
% \email{jachym.bartik@matfyz.cuni.cz}

% \author{Alžběta Šrůtková}
% \affiliation{%
%   \institution{Charles University}
%   \city{Prague} 
%   \country{Czech Republic}
% }
% \email{srutkova.alzbeta@gmail.com}

% \author{Irena Holubová}
% \orcid{0000-0003-2113-1539}
% \affiliation{%
%   \institution{Charles University}
%   \city{Prague} 
%   \country{Czech Republic}
% }
% \email{irena.holubova@matfyz.cuni.cz}

% \renewcommand{\shortauthors}{J. Bártík et al.}

% \begin{abstract}
\section*{Abstract}

The \textbf{V}ariety aspect of Big Data poses challenges for most established single-model data management solutions. However, as multi-model systems increase, especially in the data-management domain, the demand for benchmarking and comparison grows. However, real-world multi-model datasets are scarce, and existing benchmarks are limited in number and versatility.
This demo paper introduces \textit{TransforMMer}, a tool to generate pseudo-realistic multi-model data from given single-model (or multi-model) datasets. It is based on a pipeline leveraging a subset of our multi-model data management tools to automate schema inference, enable user-defined schema modifications, and generate data reflecting these changes. In the demonstration, we will introduce the tool, the repository of created multi-model datasets, their challenges, and future extensions.
% \end{abstract}

% The CCS format is removed on purpose for the camera-ready version.
%%
%% The code below is generated by the tool at http://dl.acm.org/ccs.cfm.
%% Please copy and paste the code instead of the example below.
%%
% \begin{CCSXML}
% <ccs2012>
% <concept>
% <concept_id>10003752.10010070.10010111.10010112</concept_id>
% <concept_desc>Theory of computation~Data modeling</concept_desc>
% <concept_significance>500</concept_significance>
% </concept>
% <concept>
% <concept_id>10003752.10010070.10010111.10011710</concept_id>
% <concept_desc>Theory of computation~Data structures and algorithms for data management</concept_desc>
% <concept_significance>500</concept_significance>
% </concept>
% </ccs2012>
% \end{CCSXML}

% \ccsdesc[500]{Theory of computation~Data modeling}
% \ccsdesc[500]{Theory of computation~Data structures and algorithms for data management}

% The keywords are removed on purpose for the camera-ready version.
% \keywords{Multi-model data, benchmarking, data generator, real-world data.}

% \maketitle

\section{Introduction}

The traditional relational data model has long been the standard for data representation, but Big Data has revealed its limitations. A key challenge is the \textbf{V}ariety of Big Data, which encompasses multiple types and formats that originate from diverse sources and are inherently adherent to different models. There are structured, semi-structured, and unstructured formats; order-preserving and order-ignorant models; aggregate-ignorant and aggregate-oriented systems; etc. The naturally contradictory features of the so-called \textit{multi-model data}, combining distinct models with contradictory features, complicates all aspects of data management, from modelling to querying. 

Many tools now support multi-model capabilities -- e.g., more than two-thirds of the top 50 database management systems (DBMSs)\footnote{\url{https://db-engines.com/en/ranking}} align with Gartner's decade-old prediction of support for multiple models~\citep{Gartner2015}. However, the lack of standards for combining models leaves each DBMS with proprietary solutions. Even the variety of query languages~\citep{Jiaheng2023Survey} highlights the state of multi-model data systems.

Selecting the optimal tool for a specific use case is challenging, and extensive benchmarking is essential to compare tools across all target use cases. While single-model benchmarks and data generators exist for standard models, their adaptation to the multi-model paradigm is complex. Multi-model test cases must account for numerous combinations of different models using various strategies such as embedding, cross-model references, or redundancy. And this limits the availability and versatility of truly comprehensive multi-model benchmarks.~\citep{CSUR-Jiaheng,unibench}


Let us illustrate the problem with an example. We are building a web application that depends on various types of data. When it comes to selecting a database solution, we face a dilemma. We want the reliability of a relational database (e.g., for products and orders), but we also need to be able to perform complex graph queries (e.g., about the users, their mutual connections, and interactions with our products). We can leverage a multi-model database to store such data; however, we are uncertain whether this approach is more efficient than storing everything in a single database system. Many single-model datasets are similar to our expected data, so we can use them to benchmark the performance of the single-database solution. However, there are very few multi-model datasets available, nor there are tools for generating them.


To address this challenge, we introduce \textit{TransforMMer}, a tool for generating virtually any multi-model dataset. Instead of relying on traditional value generators with predefined distributions, our approach uses real-world single-model (or multi-model) datasets as input. It infers a schema of the input data, enables its user-specified modification to reflect the required combination of models, and transforms the output to the respective form. It also supports redundancy and versioning, which further increase its versatility. For its implementation, we integrate and extend selected tools from our toolset based on the unifying categorical representation of multi-model data~\citep{DBLP:journals/jbd/KoupilH22}. Combining and extending tools for modelling and transformations~\citep{mm_cat}, schema inference~\citep{mm_infer}, and evolution management~\citep{mm_evocat} with appropriate integration modules and interfaces allows us to ensure the transformation universally and correctly.

\paragraph{Outline} In Section~2, we review related work. Section~3 introduces the categorical representation of multi-model data. Section~4 introduces \textit{TransforMMer} and Section~5 its demonstration.

%%%%%%%%%%%%%%%%%%%%%%%%%%%%%%%%%%%%%%%%%%%%%%%%%%%%%

\section{Related Work}
\pLabel{sec:rel}
Two main approaches to benchmarking data management tools involve real-world data sets or synthetic data generators. Many focus on single models, but there are very few multi-model options.

Real-world datasets are primarily relational due to the dominance of relational DBMSs, accompanied by hierarchical models (e.g., JSON for NoSQL DBMSs) and graph data. Popular sources include Kaggle\footnote{\url{https://www.kaggle.com/}} or Harvard Dataverse\footnotei{\url{https://dataverse.harvard.edu/}}{.} Graph data can be found in collections like SNAP\footnotei{\url{https://snap.stanford.edu/data/}}{.} Governments (e.g., US\footnotei{\url{https://data.gov/}}{,} EU\footnotei{\url{https://data.europa.eu/}}{,} etc.) provide open data portals involving distinct data formats. Moreover, various datasets can also be found on GitHub or Google Dataset Search.

However, often we cannot easily find a suitable real-world dataset, and we need to use a data generator or benchmark. However, most are limited to specific data models, formats, or fixed use cases. For example, TPC-H and TPC-DS\footnote{\url{https://www.tpc.org/}} focus on the relational model. Benchmarks such as XMark~\citep{10.5555/1287369.1287455} are tailored to the document model. And also, many graph data generators exist~\citep{10.1145/3379445}. 

Regarding multi-model data, BigBench~\citep{10.1145/2463676.2463712} covers semi-structured and unstructured data, but lacks support for graph and array data models. UniBench~\citep{unibench} does not support the array data model either, and it considers only a single use case. M2Bench~\citep{10.14778/3574245.3574259} encompasses relational, document, graph, and array data models but also involves only three predefined use cases.

%%%%%%%%%%%%%%%%%%%%%%%%%%%%%%%%%%%%%%%%%%%%%%%%%%%%%%%%%%%%%

\section{Categorical Model of Data}
\pLabel{sec:cat}

First, to unify the terminology from different models, we use the following terms: A \textit{kind} corresponds to a class of items (e.g., a relational table or a collection of JSON documents). A \textit{record} corresponds to one item of a kind (e.g., a table row or a JSON document). A record consists of simple or complex \textit{properties} having their \textit{domains}.

We also recall the basic notions of category theory.
A \textit{category} $\mathbf{C} = (\mathcal{O}, \mathcal{M}, \circ)$ consists of a set of objects $\mathcal{O}$, set of morphisms $\mathcal{M}$, and a composition operation $\circ$ over the morphisms ensuring transitivity and associativity. Each morphism is modelled as an arrow $f : A \rightarrow B$, where $A, B \in \mathcal{O}$, $A = dom(f), B = cod(f)$. And there is an \textit{identity} morphism $1_A \in \mathcal{M}$ for each object $A$. A category can be visualized as a multigraph, with objects represented as vertices and morphisms as directed edges.

The core and the integration point of all our tools is an abstract representation of multi-model data called \textit{schema category}~\citep{DBLP:journals/jbd/KoupilH22}. It is defined as a tuple $\mathbf{S} = (\mathcal{O}_\mathbf{S}, \mathcal{M}_\mathbf{S}, \circ_\mathbf{S})$. \textit{Objects} in $\mathcal{O}_\mathbf{S}$ correspond to the domains of the ER model's entity types, attributes, and relationship types. 
Each schema object $o \in \mathcal{O}_\mathbf{S}$ is internally represented as a tuple $(key$, $label$, $superid$, $ids)$, where $key$ is an automatically assigned internal identity, $label$ is an optional user-defined name, $superid \ne \emptyset$ is a set of attributes (each corresponding to a signature of a morphism) forming the actual data contents a given object is expected to have, and $ids \subseteq \mathcal{P}(superid)$, $ids \ne \emptyset$ is a set of particular identifiers (each modeled as a set of attributes) allowing us to distinguish individual data instances uniquely.
%
\textit{Morphisms} connect appropriate pairs of objects. Each morphism $m \in \mathcal{M}_\mathbf{S}$ is represented as a tuple $(signature$, $dom$, $cod$, $label)$. The $signature$ allows us to distinguish all morphisms except the identity ones mutually. $dom$ and $cod$ represent the domain and codomain of the morphism. Finally, $label \in \{$ \texttt{\#property}, \texttt{\#role}, \texttt{\#isa}, \texttt{\#ident} $\}$ allows us to further distinguish morphisms with semantics \uv{has a property}, \uv{has an identifier}, \uv{has a role}, or \uv{is a}. The explicitly defined morphisms are denoted as \textit{base}, obtained via the composition $\circ$ as \textit{composite}.


The decomposition of a schema category $\mathbf{S}$, eventually partial or overlapping, is defined via a set of \textit{mappings}. For each kind, the mapping specifies where and how its records are stored in a selected single-/multi-model DBMS using a so-called \textit{access path}, which recursively describes the structure of a kind.

%Sample schema categories are provided in Figs.~\pRef{figure:sc-initial}--\pRef{figure:sc-multi}. 
For example, in Fig.~\pRef{figure:sc-initial}, we can see a sample schema category with kinds \textit{User}, \textit{Business} and \textit{Tip} having root objects denoted with green. They are mapped to JSON document collections from Fig.~\pRef{figure:data-initial}. For simplicity, we omit the most common edge label \uv{has a property} and replace uninteresting properties with dots in an oval.
\section{TransforMMer}

To solve the indicated open problems, \textit{TransforMMer} enables the transformation of virtually any input dataset (single- or multi-model) to a required output multi-model dataset. 

\subsection{Sample Scenario}
Let us consider a simple scenario to demonstrate the purpose and advantages of \textit{TransforMMer}. We want to use the subset of the Yelp Open Dataset\footnotei{\url{https://www.yelp.com/dataset}}{,} namely the data that describes businesses, tips, and users, represented using the JSON format as depicted in Fig.~\pRef{figure:data-initial}. 

\begin{figure}[ht]
    \centering
    \includegraphics[width=0.8\textwidth]{papers/edbt_2025/figures/data-initial.png}
    \caption{Input Yelp dataset in the document model}
    \pLabel{figure:data-initial}
\end{figure}

We want to transform the data into a combination of the relational and graph models, e.g., the one depicted in Fig.~\pRef{figure:data-final}.

\begin{figure}[ht]
    \centering
    \includegraphics[width=0.8\textwidth]{papers/edbt_2025/figures/data-final.png}
    \caption{Yelp dataset in the graph and relational model}
    \pLabel{figure:data-final}
\end{figure}

%In order to cover real-world use cases, we continuously tested our tool on realistic datasets. We have encountered several challenges, which we have addressed by adding new features. We will demonstrate them using the Yelp dataset as an example. 

\subsubsection{Initial Schema Category}

\textit{TransforMMer} first infers the \textit{initial schema category} $\mathbf{S_{ini}}$ depicted for the input data in Fig.~\pRef{figure:sc-initial}. 
As we can see, each original JSON document corresponds to one kind in the schema category, namely \textit{User}, \textit{Business}, and \textit{Tip}.

%The dataset contains several kinds. Let's focus on \textit{business}, \textit{user}, and \textit{tip}. Our tool infers the schema of each kind, which is depicted in Fig.~\pRef{figure:sc-initial}.

\begin{figure}[ht]
    \centering
    \includegraphics[width=0.8\textwidth]{papers/edbt_2025/figures/sc-initial.png}
    \caption{Initial schema category $\mathbf{S_{ini}}$ inferred from the sample data in Fig.~\pRef{figure:data-initial}}
    \pLabel{figure:sc-initial}
\end{figure}

\subsubsection{Improved Schema Category}
As $\mathbf{S_{ini}}$ is very simple and even incorrect, \textit{TransforMMer} enables its improvements. During the process of inference of $\mathbf{S_{ini}}$, it also infers a list of candidates for basic integrity constraints (identifiers and references), clustering, or recursion. The user can confirm or refuse them or add his/her suggestions to create an improved schema category $\mathbf{S_{imp}}$.

\paragraph{Identifiers}
First, we need an identifier for each kind to create a correct schema category. The inference process of \textit{TransforMMer} creates a list of simple properties that are candidate identifiers, i.e., they have a unique and compulsory value in each respective record. The user can select the correct ones from them. The candidates do not include composite identifiers, as finding all possible combinations would be too computationally expensive. However, the user can define composite identifiers manually.

In the sample $\mathbf{S_{ini}}$ in Fig.~\pRef{figure:sc-initial}, kind \textit{Business} has only one candidate identifier \textit{business\_id}, so we choose it as its identifier. %However, if we had used only a subset of the data, there might have been multiple candidate identifiers. In that case, we must rely on the domain knowledge to select the correct one. 
Kind \textit{User} also has only one candidate identifier \textit{user\_id}. Kind \textit{Tip} does not have any (simple) candidate identifier, so we create a composite identifier consisting of properties \textit{user\_id} and \textit{business\_id}. (The identifiers are depicted in Fig.~\pRef{figure:sc-final} in curly brackets in root nodes of the kinds.)

\paragraph{References}
Next, we can define references between the kinds to the respective identifiers. \textit{TransforMMer} also automatically infers candidates for the references (based on the inclusion of values of the referenced properties).

In the sample $\mathbf{S_{ini}}$ in Fig.~\pRef{figure:sc-initial}, we identify a reference from kind \textit{Tip} to kinds \textit{User} and \textit{Business} based on the values of properties \textit{user\_id} and \textit{business\_id}. The result is depicted in Fig.~\pRef{figure:sc-final}, where the two objects \textit{user\_id} were merged into one. The same happened for objects \textit{business\_id}. However, the primary identifier of kind \textit{Tip} is still the same; it is just represented by different schema morphisms.


\paragraph{Clustering of Properties}
Further improvements in the schema category involve the integration of various more complex constructs. One of them is replacing (huge) clusters of properties with maps. 
According to our analyses of real-world data, they may often contain multiple properties which have the same structure but different names (such as the property \textit{attributes} of kind \textit{Business} depicted in Fig.~\pRef{figure:sc-initial}). They make the schema unnecessarily complex, so \textit{TransforMMer} allows the user to transform them into a single set of key-value pairs. The user must just select all the properties to be represented as a map.  

The result is depicted in Fig.~\pRef{figure:sc-final}. As we can see, the direction of morphism between objects \textit{Business} and \textit{attributes} changed. This is because initially, the object \textit{attributes} corresponded to a complex property containing all the attributes, but now, it represents an array of key-value pairs. So, the cardinality specified by the morphism changed, and the morphism had to be updated (for details, see~\citep{DBLP:journals/jbd/KoupilH22}).

\begin{figure}[ht]
    \centering
    \includegraphics[width=0.8\textwidth]{papers/edbt_2025/figures/sc-final.png}
    \caption{The improved schema category $\mathbf{S_{imp}}$ of the Yelp data containing identifiers, a reference, and a map of properties}
    \pLabel{figure:sc-final}
\end{figure}

%Some datasets also use identifiers as property keys (e.g., UNESCO\footnotei{\url{https://ich.unesco.org/en/open-access-to-dive-data-01218}}){.} This is especially unfortunate because there is an arbitrary number of them (contrary to the large, but still manageable amount of business attributes in the Yelp dataset). The tool can handle it using the previous approach, but we are working on a better solution that would allow the users to define a pattern instead of manually selecting all the keys.

\paragraph{Recursion}
Similarly to clustering of properties, i.e., generally discovering a repeating pattern, we can also identify the recursion, i.e., a pattern occurring not within the same parent property but mutually nested. For example, consider a complex property \textit{comment} consisting of properties \textit{author}, \textit{content}, and \textit{replies}, where property \textit{replies} is an array of properties \textit{comment}. Such structure can be repeated in the inferred schema as many times as is the maximum depth of recursion in the source data.

We can significantly simplify this schema by defining a pattern for the property \textit{comment}. \textit{TransforMMer} then creates only a single complex property for each unique property in the pattern and a morphism from property \textit{replies} to property \textit{comment}. Thus, the resulting schema category contains a cycle, allowing infinite recursion. An example of such a transformation is depicted in Fig.~\pRef{figure:recursion}, which simultaneously provides sample screenshots of \textit{TransforMMer}.

\begin{figure}[h!]
%    \centering
%    \begin{minipage}{0.58\textwidth}
        \centering
        \includegraphics[width=\textwidth]{papers/edbt_2025/figures/recursion-initial-novalue.png}
%    \end{minipage}%
%    \begin{minipage}{0.42\textwidth}
%        \centering
        \includegraphics[width=\textwidth]{papers/edbt_2025/figures/recursion-final-wide-novalue.png}
%    \end{minipage}
    \caption{Screenshot of \textit{TransforMMer} before and after the elimination of a recursive structure}
    \pLabel{figure:recursion}
\end{figure}

\subsubsection{Mapping Modification}
Having the final improved schema category $\mathbf{S_{imp}}$ (see Fig.~\pRef{figure:sc-final}), the user can now define its new mapping. As depicted in Fig.~\pRef{figure:sc-multi}, for example, the original document model (green) can be entirely discarded, and instead of it, the user may define four new kinds. Kind \textit{User} is mapped to the graph model (blue), where each user represents a node in the graph, and edges represent mutual friendship. Kinds \textit{Business}, \textit{Tip}, and \textit{attributes} are mapped to the relational model (purple), i.e., tables. Considering the relational model without arrays, the kind \textit{attribues} had to be established. The output data are depicted in Fig.~\pRef{figure:data-final}.

\begin{figure}[h!]
    \centering
    \includegraphics[width=0.8\textwidth]{papers/edbt_2025/figures/sc-multi.png}
    \caption{$\mathbf{S_{imp}}$ of the Yelp data re-mapped to the graph (blue) and relational (purple) model to represent data in Fig.~\pRef{figure:data-final}}
    \pLabel{figure:sc-multi}
\end{figure}

\subsection{Workflow}
The whole process of transforming the input (single- or multi-model) data to required multi-model data is depicted in Fig.~\pRef{figure:workflow}.

\begin{figure}[h!]
    \centering
    \includegraphics[width=\textwidth]{papers/edbt_2025/figures/workflow.png}
    \caption{The workflow of \textit{TransforMMer}}
    \pLabel{figure:workflow}
\end{figure}

\paragraph{Step 1}
First, the users need to specify input data sources. \textit{TransforMMer} supports the input from a file system, a DBMS, or their combination. It then infers the schema of each kind (table, collection, etc.) in the data sources. The output of this step is $\mathbf{S_{ini}}$ and its mapping to the logical structures of the original data. The algorithm also infers candidates for identifiers, references between the kinds, and repeating patterns. %An example of such SC is shown in Fig.~\pRef{figure:sc-initial}.

\paragraph{Step 2}
$\mathbf{S_{ini}}$ is improved by interacting with the user, i.e. $\mathbf{S_{imp}}$ is created. The user can specify identifiers, references, and repeating patterns to be merged from the candidate set or his/her own.
%The inferred SC can not be used directly. In order to produce multi-model data, the users need to specify identifiers of the kinds and references between them. This is done semi-automatically (see section \pRef{section:real-data}). The results are shown in Fig.~\pRef{figure:sc-final}.
The user can also further modify the schema category, e.g., by grouping a set of properties into a new property or deleting a property. With every change, \textit{TransforMMer} automatically updates the mappings. 


\paragraph{Step 3}
$\mathbf{S_{imp}}$ is mapped to the new requested combination of models. The user selects parts of the schema category and specifies the requested, possibly overlapping mapping (e.g., by modifying an existing one or creating a new one from scratch). \textit{TransforMMer} ensures all the respective modifications.

%For example, the original data was in JSON files (i.e., document model). However, we want to transform it into a combination of PostgreSQL (relational) and Neo4j (graph) databases (see Fig.~\pRef{figure:sc-multi}).

\paragraph{Step 4}
\textit{TransforMMer} loads the original data according to the input mappings into an in-memory categorical representation and then stores it according to the new user-specified mappings into the output. It can be stored in a specified file system or a DBMS.


\subsection{Architecture}
As mentioned in the introduction, \textit{TransforMMer} leverages selected members of our toolset for multi-model data management, namely \textit{MM-infer}~\citep{mm_infer} for categorical schema inference, \textit{MM-evocat}~\citep{mm_evocat} for schema modifications, and \textit{MM-cat}~\citep{mm_cat} (now a part of \textit{MM-evocat}) for schema modelling and data transformations. Fig.~\pRef{figure:workflow} also indicates their integration in \textit{TransforMMer}. \textit{TransforMMer} consists of the backend (written in Java), which contains all the algorithms, and a single-page client (written in Vue), which provides the user interface. The code repository is available at GitHub\footnotei{\url{https://github.com/mmcatdb/mmcat}}{.} And a demo is also available\footnotei{\url{https://demo.mmcatdb.com}}{.} Lastly, we have a short video tutorial\footnotei{\url{https://www.youtube.com/watch?v=Aiw815OaAb4}}{.}

All of the algorithms mentioned above use abstract wrappers, which provide a unified interface to the data sources instead of working with them directly. This allows us to easily add support for new data sources by implementing a new wrapper. For now, we support PostgreSQL, MongoDB, Neo4j (i.e., representatives of the relational, document, and graph model), CSV and JSON files (i.e., the most common file formats).

Because the users might want to work with large datasets, all processes are implemented as asynchronous jobs. %The users can start the jobs and then return later to check the results.

\subsection{Dataset Repository}
As a proof of concept, we have used \textit{TransforMMer} to generate a repository\footnote{\url{https://dare.mmcatdb.com}} of popular multi-model datasets (to be gradually extended). The current datasets involve, e.g., the Yelp dataset, the IMDb dataset\footnotei{\url{https://developer.imdb.com/non-commercial-datasets/}}{,} or the dataset of NASA open-source code projects\footnotei{\url{https://data.nasa.gov/Software/NASA-open-source-code-projects-with-A-I-generated-/3efg-u4v8/about-data}}{.} In the next steps, we want to extend the repository with the popular datasets and their variations and using a comparative survey demonstrate the wide usability of \textit{TransforMMer} for comparing multi-model databases.

%\section{The Real Data Experience}\pLabel{section:real-data}


\section{Demonstration Outline}
In our presentation, we will first demonstrate the described functionality of \textit{TransforMMer} using the Yelp dataset. We will also introduce the repository of multi-model data sets created using \textit{TransforMMer} and show selected specifics and challenges we encountered in the real-world datasets, the support for redundancy, and the applicability of versioning to simulate data evolution. Finally, we will overview future plans for the generation of respective queries. The interested audience can also experiment with \textit{TransforMMer} and own datasets.


% \section*{Acknowledgement}
% Supported by the GAČR grant no. 23-07781S, GAUK grant no. 292323., and SVV project no. 260 698.

%%
%% The next two lines define the bibliography style to be used, and
%% the bibliography file.
% \bibliographystyle{ACM-Reference-Format}
% \bibliography{bibliography}

% %%
% %% If your work has an appendix, this is the place to put it.
% %% Please note that all the content must fit within the page limits, including any appendices.
% %\appendix
% %
% %\section{Research Methods}
% % ...

% \end{document}
% \endinput
